\honest{Money: The Unit of Caring}{Eliezer Yudkowsky}{2009}

Steve Omohundro has pointed out that a rational agent should spread its resources so that shifting a unit between any two tasks gains nothing. So any such agent has a common currency of ``expected utilons.'' In our society this currency is money: ``It is the measure of how much society cares about something.'' \nb{Summed across people, spending counts each person's caring in proportion to their money. The post itself later says that money measures caring only ``up to a positive scalar factor.''} This point, I say, is ``brutal yet obvious,'' and many are ``motivated to deny'' it.

Many people who want to help would rather volunteer a few hours, mail in an old laptop, or march, than spend money. I understand: every time I spend money I feel like I am ``losing hit points.'' \nb{Behavioural economists call this the pain of paying (Prelec and Loewenstein, 1998).}

But there is ``this very, very old puzzle'' in economics: a lawyer volunteers an hour at a soup kitchen, when an extra hour of work would pay for five hours of a hired worker. \nb{No economics source for the puzzle before this post could be traced. Peter Unger had argued in 1996 that a law professor could earn far more at a corporate firm and give most of it away.} Comparative advantage and specialization are why we have money. ``This is what grownups do.'' Volunteering makes sense for people who cannot sell more of their time, or when the soup kitchen needs a large block of the lawyer's own legal work. Failing to grasp this ``may be sufficient by itself'' to prevent collective action in groups larger than 40 people. Specialization and trade ``are the only way that anything ever gets done in this world.'' \nb{Yochai Benkler (2002) described large projects built from unpaid contributions of a few minutes each, among them NASA's Clickworkers, who marked craters on Mars. Benkler's condition is that the work can be split into small pieces.}

No one raises cattle to avoid the pain of paying for beef. The idea that someone who gives ``mere money'' does not care enough to get involved is ``a mere reflection of a mind that doesn't understand'' a market economy. I give no evidence for this diagnosis. Giving money is every hour you worked to earn it, spent as if at the charity ``in your professional capacity.'' The lawyer may volunteer to stay motivated, as long as he also gives money. ``I may post more about this later.''

Money is the unit of relative caring. If you would spend \$5 on a burrito, whatever you will not spend \$5 on you care about less. Someone always reluctant to spend any money on their partner, who spends \$1000 on a flat-screen television without a qualm, has shown their relative values. \nb{Peter Singer had made a comparison of this kind about charity in 1972: people spend on new clothes or a new car ``instead of giving it to famine relief,'' and, Singer argued, ``we ought to give the money away.''} Frugality is a virtue and spending hurts, but ``if you are never willing to spend any units of caring, it means you don't care.''
